% from ega0-19-fr.tex
\oldpage[0\textsubscript{IV-1}]{86}

l'ensemble des $w_\lambda\in\operatorname{Hom}(B,C_\lambda)$ tels que
$u_\lambda$ se factorise en
$B\xrightarrow{w_\lambda}C_\lambda\to
C_\lambda/\mathfrak J_\lambda$~; $W_\lambda$ est une variété
linéaire fermée dans le $C_\lambda$-module
  $\operatorname{Hom}(B,C_\lambda)$, non vide puisque $B$ est formellement
  lisse. Dans le produit
$\prod_{\mu\leq\lambda}\operatorname{Hom}(B,C_\mu)=E_\lambda$,
considérons la sous-variété linéaire $U_\lambda$ formée des
$(w_\mu)_{\mu\leq\lambda}$ tels que
$w_\mu=\varphi_{\mu\lambda}\circ w_\lambda$ pour
$\mu\leq\lambda$ et $w_\lambda\in W_\lambda$, où
$\varphi_{\mu\lambda}:C_\lambda\to C_\mu$ est
l'homomorphisme canonique, qui est aussi fermée. Enfin, soit $V_\lambda$
la variété linéaire dans le produit
$\prod_\mu\operatorname{Hom}(B,C_\mu)$, produit de $U_\lambda$
et des $\operatorname{Hom}(B,C_\mu)$ pour
$\mu\not\leq\lambda$, qui est encore fermée. Tout revient à voir que
l'intersection des $V_\lambda$ est \emph{non vide}, car un élément $w$ de
cette intersection appartiendra à
$\varprojlim_\lambda\operatorname{Hom}(B,C_\lambda)$ par définition.
D'ailleurs, comme $\mathfrak J$ est fermé, et
$C=\varprojlim_\lambda C_\lambda$ linéairement compact,
$C/\mathfrak J$ s'identifie à
$\varprojlim_\lambda C_\lambda/\mathfrak J_\lambda$ et
l'application canonique $\psi:C\to C/\mathfrak J$ à
$\varprojlim_\lambda\psi_\lambda$, où $\psi_\lambda$ est
l'application canonique
$C_\lambda\to C_\lambda/\mathfrak J_\lambda$. On conclut donc
comme dans (19.3.10) que $\psi\circ w=u$.
\end{proof}

\subsection{Premiers critères de lissité formelle}
\label{subsection:0.19.4-fr}

\begin{proposition}[19.4.1]
\label{0.19.4.1-fr}
Soient $A$ un anneau topologique, $B$ une $A$-algèbre topologique~; supposons
qu'il existe deux familles filtrantes décroissantes
$(\mathfrak J_\alpha)_{\alpha\in I}$,
$(\mathfrak R_\alpha)_{\alpha\in I}$ d'idéaux de $A$ et $B$
respectivement, telles que~:
\begin{enumerate}
  \item[$1^\circ$] $(\mathfrak J_\alpha)$ tend vers $0$ dans $A$ et
  $(\mathfrak R_\alpha)$ tend vers $0$ dans $B$~;
  \item[$2^\circ$] pour tout $\alpha\in I$ on a
  $\mathfrak J_\alpha B\subset\mathfrak R_\alpha$ (de sorte que
  $B/\mathfrak R_\alpha$ est une
  $(A/\mathfrak J_\alpha)$-algèbre topologique)~;
  \item[$3^\circ$] pour tout $\alpha\in I$,
  $B/\mathfrak R_\alpha$ est une
  $(A/\mathfrak J_\alpha)$-algèbre formellement lisse.
\end{enumerate}
Alors $B$ est une $A$-algèbre formellement lisse.
\end{proposition}

\begin{proof}
En effet, soient $C$ une $A$-algèbre discrète, $\mathfrak L$ un idéal
nilpotent de $C$, $\mathfrak R$ un idéal ouvert de $B$~; par hypothèse il y a
un $\alpha\in I$ tel que
$\mathfrak R_\alpha\subset\mathfrak R$, donc $B/\mathfrak R$ est
quotient de $B/\mathfrak R_\alpha$ par un idéal ouvert. Tout
$A$-homomorphisme $u':B/\mathfrak R\to C/\mathfrak L$ est aussi un
$(A/\mathfrak J_\alpha)$-homomorphisme, donc il existe un idéal ouvert
$\mathfrak R'$ de $B$ tel que
$\mathfrak R_\alpha\subset\mathfrak R'\subset\mathfrak R$ et que
$B/\mathfrak R_\alpha\to B/\mathfrak R\xrightarrow{u'}C/\mathfrak L$
se factorise en
$B/\mathfrak R_\alpha\to B/\mathfrak R'\to C\to C/\mathfrak L$,
où le second homomorphisme est un
$(A/\mathfrak J_\alpha)$-homomorphisme~; d'où la conclusion, en vertu de
la remarque suivant (19.3.1).
\end{proof}

\begin{corollary}[19.4.2]
\label{0.19.4.2-fr}
Soient $A$ un anneau topologique, $B$ une $A$-algèbre topologique,
$(\mathfrak J_\alpha)_{\alpha\in I}$ une famille filtrante décroissante
d'idéaux de $A$ tendant vers $0$. Pour que $B$ soit une $A$-algèbre
formellement lisse, il faut et il suffit que pour tout $\alpha\in I$, si
l'on pose $A_\alpha=A/\mathfrak J_\alpha$,
$B_\alpha=B/\mathfrak J_\alpha B$, $B_\alpha$ soit une
$A_\alpha$-algèbre formellement lisse.
\end{corollary}

La condition est suffisante par (19.4.1), et elle est aussi nécessaire par
(19.3.5), (iii).

\begin{proposition}[19.4.3]
\label{0.19.4.3-fr}
Soient $A$ un anneau topologique, $B$ une $A$-algèbre topologique. Supposons
que pour toute $A$-algèbre discrète $C$ et tout idéal $\mathfrak J$ de $C$
tel que $\mathfrak J^2=0$, tout $A$-homomorphisme continu
$u:B\to C/\mathfrak J$ se factorise en
$B\xrightarrow{v}C\to C/\mathfrak J$, où $v$ est un $A$-homomorphisme
continu. Alors $B$ est une $A$-algèbre formellement lisse.
\end{proposition}

\begin{proof}
En effet, avec les mêmes notations, soit $\mathfrak R$ un idéal nilpotent
quelconque de $C$, et considérons un $A$-homomorphisme continu
$u':B\to C/\mathfrak R$. Supposons que $\mathfrak R^m=0$, et posons
$C_i=C/\mathfrak R^i$ pour $1\leq i\leq m$, de sorte que $C_m=C$, que
$\mathfrak R^{i-1}/\mathfrak R^i$ est un idéal de carré nul
$\mathfrak J_i$ dans $C_i$, et $C_i=C_{i+1}/\mathfrak J_{i+1}$~;
l'hypothèse entraîne alors par récurrence sur $i$ l'existence de
$A$-homomorphismes continus $u_i:B\to C_i$ tels que $u_1=u$ et que $u_i$
se factorise en
$B\xrightarrow{u_{i+1}}C_{i+1}\to C_{i+1}/\mathfrak J_{i+1}=C_i$~;
d'où la conclusion.
\end{proof}

\begin{proposition}[19.4.4]
\label{0.19.4.4-fr}
Soient $A$ un anneau topologique, $B$ une $A$-algèbre topologique
(commutative). Pour que $B$ soit une $A$-algèbre formellement lisse, il faut et
il suffit que pour

